\section*{Chapter \ref{ch:b3:demographic-transition} --- Population and the Demographic Transition}

\begin{solution}{exo:b3:demographic-transition:1}
$\ln(1{,}611/1{,}179)/350 = 0.09~\%$ a year. After 1650 income per person rises by half in fifty years, and
from 1750 it grows by more than 0.35~\% a year in every half-century. The single years
mislead because the series swings by 10 to 15~\% from year to year: 1650 is a trough
(1{,}446, against 1{,}683 in 1640 and 1{,}982 in 1655), which gives 0.06~\% from 1300 to
1650 and exaggerates the rise after it. % figdata: england_stats
\end{solution}

\begin{solution}{exo:b3:demographic-transition:2}
Population is constant at $y = 400$, so $1{,}000 - 0.5L = 400$ and $L = 1{,}200$ thousand.
\end{solution}

\begin{solution}{exo:b3:demographic-transition:3}
$0.45 \times 0.40 = 18~\%$ of wages.
\end{solution}

\begin{solution}{exo:b3:demographic-transition:4}
Income per person rising for a century or more without a rise in population;
better land or technology going with higher living standards and not with denser
population; or wages that did not rise after large falls in population.
\end{solution}

\begin{solution}{exo:b3:demographic-transition:5}
With $y = A L^{-1/2}$, income doubles at once; population then grows, faster at first, and
income falls back towards $\bar y$. In the long run $L^* = (A/\bar y)^2$ is four times larger and income is unchanged.
\end{solution}

\begin{solution}{exo:b3:demographic-transition:6}
$5 \times (0.02 + 0.10 + 0.12 + 0.09 + 0.04 + 0.01) = 5 \times 0.38 = 1.9$. The \omterm{def:b3:demographic-transition:tfr}{net reproduction rate} is $1.9 \times 0.488 \times 0.98 = 0.91$: below replacement.
\end{solution}

\begin{solution}{exo:b3:demographic-transition:7}
The tested solution is \code{ex07_fertility_income.py}, in the chapter's
\code{exercises/} folder. In 1960, 107 countries, mean 5.51 children per woman and a slope of $-2.04$ per
tenfold income; in 2023, 203 countries, mean 2.37 and slope $-1.49$. Fertility fell at every level
of income, more where it was high. % output: ex07_fertility_income
\end{solution}

\begin{solution}{exo:b3:demographic-transition:8}
$n$ is proportional to $a - b$: from 0.3 to 0.2, a factor $2/3$. $q$ is proportional to $b/(a - b)$: from $0.667$ to 1.5, a factor 2.25.
\end{solution}

\begin{solution}{exo:b3:demographic-transition:9}
If an unplanned twin birth leaves the schooling of the older children unchanged,
the trade-off fails at that margin, as in Norway and Israel. A fall in their schooling,
beyond what the twins' low birth weight explains, supports it, as in China.
\end{solution}

\begin{solution}{exo:b3:demographic-transition:10}
$d\ln L/dt = h(A L^{-(1 - \alpha)})$. Its derivative with respect to $\ln L$ at $L^*$ is $-(1 - \alpha)\bar y\,h'(\bar y) < 0$, so $\ln L$ returns to $\ln L^*$ from
either side, at the speed $(1 - \alpha)\bar y\,h'(\bar y)$. With $h(y) = \phi(\ln y - \ln \bar y)$ this is $\phi(1 - \alpha)$, as in the text.
\end{solution}

\begin{solution}{exo:b3:demographic-transition:11}
The tested solution is \code{ex11_half_centuries.py}, in the chapter's
\code{exercises/} folder. The rates are $+0.26$, $+0.52$, $-0.09$, $+0.03$, $-0.01$, $-0.05$, $-0.04$, $+0.85$, $+0.17$, $+0.38$ and $+0.58~\%$ a year; from 1750 growth stays
above 0.3~\%. % output: ex11_half_centuries
The 1350--1400 rise is the Black Death's; the 1650--1700 rise is the largest before 1750.
\end{solution}

\begin{solution}{exo:b3:demographic-transition:12}
Log growth is $0.40 + \ln 1.2 = 0.40 + 0.18 = 0.58$ over twenty years, 2.9~\% a year; the dividend is $0.18/0.58 = 31~\%$ of it.
\end{solution}

\begin{solution}{pb:b3:demographic-transition:factories}
% figdata: agrarian_stats
\textbf{1.} Population must grow at $2 \times 0.1 = 0.2~\%$ a year, so $0.1\ln(y/300) = 0.002$: $y = 300e^{0.02} = 306.1$ kilograms.

\textbf{2.} Population grows only when income exceeds 300, and it must grow to keep pace with technology.

\textbf{3.} $(9{,}000/306.1)^2 = 864.7$ thousand.

\textbf{4.} 1{,}054.1 thousand, growing 0.2~\% a year.

\textbf{5.} Income would return to 306; population would be four times larger.

\textbf{6.} No trend: the real wage stays at 306 kilograms.

\textbf{7.} 633.7 thousand.

\textbf{8.} 395.1 kilograms, a premium of 29.1~\%.

\textbf{9.} $1/\sqrt{0.6} = 1.291$: the same.

\textbf{10.} At the maximum, 1.2~\% a year: the formula gives $0.1\ln(395.1/300) = 2.75~\%$, above the cap.

\textbf{11.} 1626.

\textbf{12.} 307.0 kilograms and 1{,}282.5 thousand: back at the steady state. England's income rose 24~\% from 1348 to 1351 and
stayed high for a century, longer than here.

\textbf{13.} 306.1 kilograms and 1{,}513.7 thousand in 1780; 587.8 and 3{,}351.9 thousand in 1850.

\textbf{14.} 1853.

\textbf{15.} 1.2~\% in 1850 and 0.72~\% in 1899: with income above 600, the maximum has started to fall --- the
fertility decline.

\textbf{16.} 969.6 kilograms and 5{,}520.5 thousand.

\textbf{17.} $n = 0.3/(1.5 \times 0.04) = 5$ and $q = 0.2 \times 0.04/(0.3 \times 0.01) = 2.67$ years of schooling.

\textbf{18.} $n = 0.15/0.06 = 2.5$ and $q = 0.014/0.0015 = 9.33$.

\textbf{19.} No: $n$ rises with $I$. Fertility falls because the weight on quality rises.

\textbf{20.} $1{,}950/4{,}366 = 0.447$.

\textbf{21.} $0.40 \times 0.447 = 17.9~\%$.

\textbf{22.} $0.40 \times 2{,}400/4{,}366 = 22.0~\%$.

\textbf{23.} $0.40 \times 1{,}950/2{,}400 = 32.5~\%$.

\textbf{24.} Older people who work pay contributions instead of drawing pensions: the
numerator falls and the denominator rises.

\textbf{25.} The steady-state income is 306 kilograms of grain a year, the plague's premium 29.1~\%, and the 2025
contribution rate 17.9~\% of wages.
\end{solution}
